\chapter*{Anexo: Glosario}

En esta sección se van a recoger las definiciones de términos que hayan sido utilizados en la memoria:

\begin{description}
    \item[\hypertarget{api}{Application Programming Interface (API)}]: conjunto de reglas, protocolos y herramientas 
    que permite que dos aplicaciones de software se comuniquen e intercambien datos entre sí de forma segura.
    
    \item[\hypertarget{singlePageApplication}{Single Page Application}]: es un tipo de aplicación web que ejecuta todo 
    su contenido en una sola página, cargando el contenido HTML, CSS y JavaScript por completo al abrir la web. Al ir
    pasando de una sección a otra, solo necesita cargar el contenido nuevo de forma dinámica si este lo requiere, 
    pero no hace falta cargar la página por completo.

    \item[\hypertarget{progressiveWebApps}{Progressive Web Apps}]: es una aplicación web que, gracias a tecnologías 
    modernas (HTML, JavaScript, CSS), ofrece una experiencia similar a una aplicación nativa, permitiendo su 
    instalación en la pantalla de inicio, funcionar sin conexión y enviar notificaciones push, todo desde una única 
    base de código y sin necesidad de tiendas de aplicaciones, combinando lo mejor de la web y las apps nativas. 

    \item[\hypertarget{clientSideRendering}{Client-Side Rendering}]: técnica que permite al navegador del usuario 
    descargar un HTML mínimo y luego usar JavaScript para construir y renderizar el contenido de la página web 
    dinámicamente en el cliente.

    \item[\hypertarget{marcosTrabajo}{Marcos de trabajo (\textit{frameworks})}]: conjunto estandarizado de herramientas, 
    librerías, reglas y componentes reutilizables que actúa como esqueleto para desarrollar software, aplicaciones web 
    o móviles de forma rápida y eficiente.

    \item[\hypertarget{jwt}{JSON Web Token}]: es un estándar abierto para la transmisión segura de la información entre dos partes 
    (típicamente un cliente y un servidor). Cada JWT es firmado digitalmente para prevenir su manipulación y además, contiene 
    \textit{claims} (piezas de información) sobre el usuario o la sesión.

    \item[\hypertarget{singleSignOn}{Single Sign-on}]: procedimiento de autenticación que habilita a un usuario determinado para 
    acceder a varios sistemas con una sola instancia de identificación.

    \item[\hypertarget{jsonb}{JSONB}]: formato binario para guardar JSON que permite indexar los datos.
\end{description}